\clearpage
\SourceLeaf{177}\begin{center}
{\Large\textls{REVIEW.}}
\end{center}

\bigskip
\hrule
\begin{center}
Göttingische gelehrte Anzeigen. \quad 1816 April 20.
\end{center}
\hrule

\bigskip
\begin{center}
{\large\textls{Stuttgart.}}
\end{center}

Printed by J. F. Steinkopf: \emph{Commentatio in primum elementorum Euclidis librum,
qua veritatem geometriae principiis ontologicis niti evincitur, omnesque
  propositiones, axiomatum geometricorum loco habitae, demonstrantur. Auctore}
  J. C. \textsc{Schwab}, \emph{Regi Württembergiae a consiliis aulicis
secretioribus, academiae scientiarum Petropolitanae, Berolinensis et Harlemensis
Sodali.} 1814. 65 pages in octavo.

\begin{center}
{\large\textls{Mainz.}}
\end{center}

Published at the author's expense and sold on commission by Florian Kupferberg:
\emph{Vollständige Theorie der Parallellinien. Nebst einem Anhange, in welchem
der erste Grundsatz zur Technik der geraden Linie angegeben wird. Herausgegeben
von} \textsc{Matthias Metternich}, \emph{Doctor der Philosophie, Professor der
Mathematik, Mitglied der gelehrten Gesellschaft nützlicher Wissenschaften zu
Erfurt.} 1815. 44 pages in octavo.

There can be few subjects in mathematics on which as much has been
written as on the gap at the beginning of geometry in establishing the
theory of parallel lines. Seldom does a year pass without some new
attempt to fill this gap coming to light, yet if we wish to speak
honestly and openly, we cannot say that in substance we have advanced any
further than Euclid had 2000 years ago. Such an
honest and forthright admission
